% sec01_1.tex — Section 1.1: Introduction

\section{Introduction}
\label{sec:1.1}

We wish to discuss the solution of elementary problems involving partial differential equations, the kinds of problems that arise in various fields of science and engineering. A \textbf{partial differential equation} (PDE) is a mathematical equation containing partial derivatives, for example,
\begin{equation}
\pd{u}{t} + 3\pd{u}{x} = 0.
\label{eq:1.1.1}
\end{equation}
We could begin our study by determining what functions $u(x,t)$ satisfy \eqref{eq:1.1.1}. However, we prefer to start by investigating a physical problem. We do this for two reasons. First, our mathematical techniques probably will be of greater interest to you when it becomes clear that these methods analyze physical problems. Second, we will actually find that physical considerations will motivate many of our mathematical developments.

Many diverse subject areas in engineering and the physical sciences are dominated by the study of partial differential equations. No list could be all-inclusive. However, the following examples should give you a feeling for the type of areas that are highly dependent on the study of partial differential equations: acoustics, aerodynamics, elasticity, electrodynamics, fluid dynamics, geophysics (seismic wave propagation), heat transfer, meteorology, oceanography, optics, petroleum engineering, plasma physics (ionized liquids and gases), and quantum mechanics.

We will follow a certain philosophy of applied mathematics in which the analysis of a problem will have three stages:
\begin{enumerate}
\item Formulation
\item Solution
\item Interpretation
\end{enumerate}

We begin by formulating the equations of heat flow describing the transfer of thermal energy. Heat energy is caused by the agitation of molecular matter. Two basic processes take place in order for thermal energy to move: conduction and convection. \textbf{Conduction} results from the collisions of neighboring molecules in which the kinetic energy of vibration of one molecule is transferred to its nearest neighbor. Thermal energy is thus spread by conduction even if the molecules themselves do not move their location appreciably. In addition, if a vibrating molecule moves from one region to another, it takes its thermal energy with it. This type of movement of thermal energy is called \textbf{convection}. In order to begin our study with relatively simple problems, we will study heat flow only in cases in which the conduction of heat energy is much more significant than its convection. We will thus think of heat flow primarily in the case of solids, although heat transfer in fluids (liquids and gases) is also primarily by conduction if the fluid velocity is sufficiently small.
